
\section{What is the \nn{}}
\label{sec:what_is}

While this concept may relate to other areas of computing, the scope will be narrowed down to kubernetes, and more so operators.
Narrowing the scope down will allow explaining the concept in a  cleaner process as there is a number of moving parts that all plays a part.
Firstly, we need a common understanding of what an operator is.

\subsection{What is an Operator}

An operator is an application that is deployed onto a cluster that will manage operands.
In the management of the operands the operator will speak to the cluster via the API server.
The API server will send events for resources on the cluster that operator has registered interest in.
These resources can be the operators self defined CRDs, third party CRDs, or kubernetes native resources.

When the operator receives events on these resources from the API server, the operator triggers a reconcile function that acts on the resource.

An operator can run in two modes, cluster scoped and namespace scoped.
The operator is far more powerfully when running in cluster scoped.
In the past few years there has been a move to have all operator be cluster scoped.
It is when the operator is running cluster scoped the problem of the \nn{} becomes easier to introduce.

\subsection{Communication between an operator and the API server}
The API server is the keeper of all information on the cluster, and an operator needs permission to speak with the API server.
This permission system is managed by RBAC allowing CRUD operations on the cluster.
So far there is nothing out of the normal about this setup.
One thing to note here is the "R" in the CRUD allows \textit{read}, \textit{list}, and \textit{watch}.

The \textit{watch} is a key part of the problem.

\subsection{So now, what is the \nn{}}
The \nn{} is an operator that knows too much about what is going on in the cluster.
When the operator set up the watch for resources the API server will send the events for of those resources.
Even if the operator only cares about a small sub set of the resources.

The RBAC rules allow the operator to receive all the events, and the reconcile functions predicates ensure only the resources the operator cares about are acted upon.
But without the cache filtering configuring of the controller-runtime, the operator will store all events in memory just in case it is needed later.
It is this storing of events just in case the data is needed, is what makes the operator a \nn{}.




